% !TeX root = ../深度学习笔记.tex
% 由主文档通过 \input 载入；请编译主文档。
\setcounter{chapter}{0} % 保留原书章号。
\chapter{深度学习的问题、方法与学习路线}
\label{chap:dl-introduction}
一张灰度图像可以表示为像素数组，但“这张图像是否包含猫”并不是逐个像素求和就能回答的问题。学习系统利用已有图像及其标签，调整一个带参数的函数，使它在未见图像上也能给出有用预测。深度学习的核心是让特征变换本身参与学习：参数既决定如何提取表示，也决定如何利用表示完成任务。

\section{从具体任务到学习问题}
设房屋面积为 $x\in\mathbb R$、价格为 $y\in\mathbb R$，由面积预测价格是\emph{回归（regression）}；由图像 $\mathbf x$ 预测有限类别 $y\in\{1,\ldots,K\}$ 是\emph{分类（classification）}；由已出现的文字预测后续词元则是序列预测。同一种输入可以服务于不同任务，任务的区别必须由目标和使用条件说明。房价还受地段等因素影响，输入中没有的信息不能仅靠增加网络层数凭空恢复。

\emph{监督学习（supervised learning）}使用输入与目标配对的数据，目标可以来自人工标注、测量或已有记录；\emph{无监督学习（unsupervised learning）}研究未给定任务标签的数据结构或分布；\emph{自监督学习（self-supervised learning）}从数据自身构造预测信号，例如遮住句中的一个词再预测它。自监督学习通常归入广义无监督学习，也可采用监督式损失训练。它们划分的是学习信号，而不是网络结构。

\section{模型、目标函数与学习算法}
\begin{definition}[模型与学习问题]
给定输入空间 $\mathcal X$、目标空间 $\mathcal Y$ 和参数集合 $\Theta$，模型族是 $\{f_\theta:\mathcal X\to\mathcal A:\theta\in\Theta\}$，其中 $\mathcal A$ 是预测空间，可以是数值、分数或概率向量。损失 $\ell:\mathcal A\times\mathcal Y\to\mathbb R\cup\{+\infty\}$ 衡量预测与目标的差异，通常越小越好；常见分类、回归损失非负，连续密度的负对数却未必非负。训练算法根据数据和规定的目标选择参数，涉及期望时另要求相应可积性。
\end{definition}
例如 $f_{w,b}(x)=wx+b$ 是模型，$\ell(\widehat y,y)=(\widehat y-y)^2/2$ 是平方损失，而梯度下降是一种求参数的方法。模型不等于损失，损失不等于算法：同一模型可以最小化平方误差或绝对误差，同一目标也可用不同算法求解。参数 $w,b$ 由训练调整，层数和学习率等\emph{超参数（hyperparameter）}通常由外部配置或验证过程选择。

\emph{训练（training）}输入数据和初始状态，输出拟合后的参数；\emph{推断（inference）}使用已经确定的参数对新输入计算结果。训练时还可能维护动量等优化器状态，它们与模型参数不同。分类模型输出一个概率向量后，仍需通过取最大值或代价相关规则形成决策。

\section{从人工特征到层次表示}
\emph{特征（feature）}是输入中供模型使用的量，\emph{表示（representation）}是数据经变换后用于后续计算的形式。若价格随面积非线性变化，可人工构造 $(x,x^2)$，再对这两个特征使用线性模型；神经网络则把特征写成可学习映射 $\mathbf h_\phi(\mathbf x)$，并联合学习预测器 $g_\psi$，形成 $f_\theta=g_\psi\circ\mathbf h_\phi$，其中 $\theta=(\phi,\psi)$。

多层组合可逐步改变数据的可分性，但不能据此认定每一层必然对应某种固定人类语义。网络深度是变换的层数，宽度是中间表示的维度，参数量则是独立可训练数值的个数；三者相关但不同。\emph{归纳偏置（inductive bias）}是模型与算法偏好的结构，例如卷积共享不同位置的局部处理规则。偏置与任务相符时可能提高数据利用效率，是否有效仍需验证。

\section{数据划分与学习结果的含义}
训练集用于拟合参数，验证集用于选择超参数与检查点，测试集用于方案确定后的独立评价。\emph{泛化（generalization）}关注模型对未参与拟合和选择的新样本的表现。记住训练样本只能证明模型能适应该数据，不能证明它已掌握可迁移规律；训练误差很低而新样本误差很高称为\emph{过拟合（overfitting）}。

划分单位必须与使用问题一致。同一人的多张照片若分别进入训练集和测试集，模型可能利用身份线索，不能据此估计对新人的表现；预测未来时，用未来数据估计标准化参数也可能引入\emph{数据泄漏（data leakage）}。因此应先确定测试对象是新样本、新对象还是未来时段，再设计划分。总体使用风险还取决于部署分布与错误代价，不能由一个测试准确率完整描述。

\section{例题与习题}
\begin{example}[高准确率的简单基线]
某测试集中有 $950$ 张非猫图像和 $50$ 张猫图像。始终预测“非猫”的准确率是多少，它能否完成寻找猫图像的任务？
\end{example}
\begin{solve}
该模型预测正确 $950$ 次，准确率为 $95\%$，但找出的猫图像数为零。它为类别不平衡任务提供了一个必须超越的简单基线，也说明总体正确比例与指定类别的识别能力不同；后者需用第\ref{chap:dl-baselines}章的召回率等指标描述。
\end{solve}
\begin{enumerate}
\item 将“预测学生是否通过考试”写成输入、目标、预测时点和评价对象明确的问题，列举一个可能泄漏结果的特征。
\item 解释为什么更换优化器没有改变模型族，而把输入从 $x$ 改为 $(x,x^2)$ 改变了可表示的函数。
\item 为存在同一对象重复观测的数据设计训练、验证和测试划分，并说明它估计哪一种泛化能力。
\end{enumerate}

